\begin{table}[!htbp]
\centering
\begingroup\singlespacing\fontsize{9.2}{11.32}\selectfont
\setlength{\tabcolsep}{3pt}
\caption{各分析分区的结局描述统计}\label{tab:descriptives}
\begin{tabular}{@{}>{\raggedright\arraybackslash}p{0.14\linewidth}>{\raggedright\arraybackslash}p{0.25\linewidth}>{\raggedright\arraybackslash}p{0.25\linewidth}>{\raggedright\arraybackslash}p{0.25\linewidth}@{}}
\toprule
结局 & 全部样本 & 开发集 & 留出集 \\
\midrule
NA & 0.367 (0.160) & 0.364 (0.159) & 0.388 (0.169) \\
PA & 0.543 (0.156) & 0.545 (0.154) & 0.524 (0.169) \\
LS & 0.419 (0.210) & 0.421 (0.211) & 0.401 (0.205) \\
\bottomrule
\end{tabular}
\tabnote{NA 为消极情绪，PA 为积极情绪，LS 为生活满意度。数值为存储单位下的均值（样本标准差）。全部样本 $n=1,037$，开发集 $n=933$，留出集 $n=104$。完整保留样本中各结局的范围为 0--1；留出集不必覆盖完整范围。}
\endgroup
\end{table}
